\subsection{The fixed joint population family}
\label{supp:inference121}

The following result specializes the
classical bounded-mean and binomial
arguments to the declared R121 family
\citep{hoeffding1963,clopperpearson1934}.
It supplies no assertion that a
particular observed comparison has
a favorable sign.

\begin{proposition}[Joint inference with interval-valued economic outcomes]
\label{prop:jointinference121}
Condition on all frozen actors,
acquisition protocols and total
statewise services.  Suppose the
64 initial-state rows are independent
and uniform on the declared product
of two 40-bit state grids, under the
stated ideal independent-entropy
assumption.  The same rows may be
used across every risk, family,
loss and budget.
Assume every retained interval
contains its indicated original
economic quantity.

There are two original-optimum
means, 58 completed-policy regret
means, and 28 signed
VALUE-minus-MARGINAL gain means.
Use interval sample means enlarged
by the following radii and
intersected with their known supports:
\[
 r_{\rm opt}=7K/32,\qquad
 r_{\rm regret}=35\tau/128,\qquad
 r_{\rm gain}=35\tau/64 .
\]
For each of the 28 gains, at each
threshold \(2^{-20},2^{-16}\),
count only certified benefit events
or certified harm events, giving
112 one-sided probability bounds.
Use lower Clopper--Pearson inversion
with error \(1/11200\) for these
counts.  For each of the 58 first
rejection indicators use lower
inversion at its count and one
minus lower inversion at its
complement count, each at error
\(1/11600\).
Forty-eight rational bisections
may replace an exact inverse by
its lower bracket.

All these statements hold jointly
with probability at least \(19/20\).
Adding each simultaneous regret
interval to its risk level's
simultaneous optimum interval
gives a population policy-objective
interval, which may be intersected
with \([0,K+\tau]\).
Runtime repetitions and risk
quadrature nodes are not additional
independent state units.
\end{proposition}

\begin{proof}
The all-state service result supplies
regret support \([0,\tau]\) and
signed gain support \([-\tau,\tau]\).
The primitive range supplies
optimum support \([0,K]\).
These are properties of the fixed
services and original objectives,
not conclusions inferred from the
64 observed successes.

For completeness, if \(Y\) is in
\([a,b]\), the centered log moment
generating function has second
derivative equal to a variance
under an exponentially tilted law.
That variance is at most
\((b-a)^2/4\), since the squared
distance from the interval midpoint
is at most this number.
Its value and first derivative at
zero are zero.  Integrating twice
gives
\[
 \log\mathbb E e^{\lambda(Y-\mathbb EY)}
                         \leq\lambda^2(b-a)^2/8.
\]
Independence across the 64 state
rows, the exponential Markov
inequality, and minimization over
\(\lambda\) yield
\[
 \mathbb P\{|\overline Y-\mathbb EY|>r\}
               \leq2\exp\{-2\cdot64r^2/(b-a)^2\}.
\]
This applies to the latent exact
outcomes.  If their sample mean
is enclosed between the means of
the reported lower and upper
endpoints, enlarging those two
endpoints by \(r\) preserves
coverage.  The numerical
enclosures themselves need not
be independent for this argument.

For regret and signed gain the
exponent is \(1225/128\); for the
original optimum it is \(49/8\).
No rounded logarithm is required.
The positive degree-20 Taylor
polynomial of \(e^x\) is a lower
bound for \(x>0\).
For \(x=1225/128\), the exact
integer floors of its terms
from degree 0 to 20 are
\[
 \begin{gathered}
 1,9,45,146,349,669,1067,1458,1745,1855,\\
 1776,1545,1232,907,620,395,236,133,70,35,17,
 \end{gathered}
\]
whose sum is \(14310>11600\).
For \(x=49/8\) the corresponding
floors are
\[
 1,6,18,38,58,71,73,64,49,33,20,11,5,2,1,
 0,0,0,0,0,0,
\]
whose sum is \(450>400\).
Each floor is verified by a
finite integer product comparison.
Therefore the two-sided regret
failure is less than \(1/5800\),
the two-sided signed-gain failure
is less than \(1/2800\), and the
two-sided optimum failure is
less than \(1/200\).

Here is the binomial argument
including numerical event
uncertainty.  For \(1\leq k\leq64\),
let
\[
 B_p(k)=\sum_{j=k}^{64}{64\choose j}
                               p^j(1-p)^{64-j}.
\]
It increases in \(p\).
The lower binomial inverse at
error \(d\) is the value
\(p_d(k)\) solving \(B_p(k)=d\);
put \(p_d(0)=0\).
If \(K_0\) is an exact binomial
count with true success
probability \(p\), the event
\(p_d(K_0)>p\) lies in an
upper-tail event of probability
at most \(d\).
To check this directly, its
smallest possible count \(k_0\)
satisfies \(B_p(k_0)\leq d\),
and the failure event is
contained in \(\{K_0\geq k_0\}\).
Replacing the inverse by a
checked lower bracket preserves
the bound.
For rational trial probabilities,
\(B_p(k)\leq d\) is an integer
comparison after clearing
denominators, so the 48
bisection steps need no
floating tail-probability
assumption.

A certified benefit count
counts rows whose gain lower
endpoint is at least the
positive threshold.
It is bounded by the latent
count of true benefits.
The harm count, using a gain
upper endpoint at most the
negative threshold, has the
same domination property.
The lower binomial inverse is
nondecreasing in the count:
increasing \(k\) decreases
\(B_p(k)\) at each \(p\).
The fixed dyadic lower
bisection bracket preserves
that ordering.
Consequently substituting a
certified count no larger
than the latent one can only
lower the probability bound.
This remains valid if the
effort used to construct
valid numerical intervals
is adaptive; the latent
state outcomes, not the
enclosure algorithm, supply
the binomial count.

First rejection is a
specified statewise binary
event for a frozen complete
service, including its rule
for a recoverable numerical
proposal failure.
Its 64 observations have
the binomial law directly.
Apply the same inverse to
its count and to the
complement count to obtain
the two one-sided endpoints.
No independence between
those endpoints, the
threshold events, or other
methods is needed.

Finally the declared union
allocation is exactly
\[
 \frac2{200}+\frac{58}{5800}
 +\frac{28}{2800}+\frac{112}{11200}
 +\frac{116}{11600}=\frac1{20}.
\]
Apply a union bound.
On this joint event the
identity
\(\mathbb EJ=\mathbb EV_0+
\mathbb E(J-V_0)\)
gives the stated policy
objective interval.
It involves the same
original state law in both
summands.
The rather wide optimum
radius is retained; a
sample of 64 states is
not turned into a precise
population level by
precise arithmetic.
\end{proof}

\subsection{Interpretation and the remaining cost obligation}

Every mathematical implication above holds
for the original economic state where its
inequalities are established.  Population
statements require their separately specified
sampling and inference assumptions.
Finite training labels alone do not establish
uniform coverage, and a uniform conservative
bound does not demonstrate that its
construction is affordable.
The complete cost comparison must pay for
the original-\(Q\) value bundles, every
probe, basis construction, bounded value
fit, marginal LP, exact primal--dual
verification, feature ranges, action
extraction, rejected certification and
the complete accepted policy service.
Any scientific reference calculation is
identified separately and cannot be
presented as free deployment work.

The derivative target eliminates neither
the original tail integration nor the
unchanged-future floor in
\eqref{eq:future121}.
It makes their placement explicit:
the original risk integration is paid
before fitting, the action-relevant
training objective is minimized over
a stated finite class, and the future
allowance is applied to an implemented
policy after the original-\(Q\) calculation.
The matched hinge, polynomial, spline,
grid and semi-analytic methods have
the same access to this construction.
A conclusion about neural adoption,
scalability or an applied economic
calibration therefore still requires
its corresponding complete-cost and
economic evidence.
